\subsection{李代数积的包络代数}

设 \(g_{1}, g_{2}\) 是 K 上的两个李代数，\(U_{i}\) 是 \(g_{i}\) 的包络代数，\(\sigma_{i}\) 是从 \(g_{i}\) 到 \(U_{i}\) 的典范映射（i = 1, 2）。令 \(g = g_{1} \times g_{2}\)，U 为其包络代数，\(\sigma\) 为从 g 到 U 的典范映射。\(g_{1}\) 和 \(g_{2}\) 到 g 的典范嵌入定义了从 \(U_{1}\) 和 \(U_{2}\) 到 U 的典范同态；它们的像彼此可交换，因而得到同态 \(\phi\)，它从代数 \(U_1 \otimes_K U_2\) 到代数 \(U\)，并将 1 映到 1。

命题 2. 同态 \(\phi\) 是代数同构

映射 \(\sigma': (x_1, x_2) \mapsto \sigma_1(x_1) \otimes 1 + 1 \otimes \sigma_2(x_2) (x_1 \in \mathfrak{g}_1, x_2 \in \mathfrak{g}_2)\) 是一个 \(\alpha\)-映射，从 \(\mathfrak{g}\) 到 \(\mathrm{U}_1 \otimes_{\mathbb{K}} \mathrm{U}_2\)，因此存在唯一的同态 \(\tau\)，它从 \(\mathrm{U}\) 到 \(\mathrm{U}_1 \otimes_{\mathbb{K}} \mathrm{U}_2\)，将 1 映到 1（第 1 号，命题 1），使得：

\[
\sigma^ {\prime} = \tau \circ \sigma .\tag{1}
\]

现在 \(\phi \circ \tau \circ \sigma = \phi \circ \sigma' = \sigma\)，且 \(\tau \circ \phi \circ \sigma' = \tau \circ \sigma = \sigma'\)，所以 \(\phi \circ \tau\) 和 \(\tau \circ \phi\) 分别是 U 与 \(U_1 \otimes_K U_2\) 的恒等映射。故命题得证。

\(U_{1} \otimes_{K} U_{2}\) 在同构 \(\phi\) 下与 U 相同。于是 g 到 U 的典范映射由 (1) 与下列映射相同：

\[
(x _ {1}, x _ {2}) \mapsto \sigma_ {1} (x _ {1}) \otimes 1 + 1 \otimes \sigma_ {2} (x _ {2}).
\]

类似地，若 \(g_{1},\ldots,g_{n}\) 是 K 上的李代数，其包络代数为 \(U_{1},\ldots,U_{n}\)，则 \(g_{1}\times\cdots\times g_{n}\) 的包络代数 U 自然地与 \(U_{1}\otimes_{K}\cdots\otimes_{K}U_{n}\) 相同，而 \(g_{1}\times\cdots\times g_{n}\) 到 U 的典范映射与下列映射相同：

\[
(x _ {1}, \dots , x _ {n}) \mapsto \sigma_ {1} (x _ {1}) \otimes 1 \otimes \dots \otimes 1 + \dots + 1 \otimes \dots \otimes 1 \otimes \sigma_ {n} (x _ {n})
\]

（\((\sigma_{i} \text{ 所表示的典范映射从 } g_{i} \text{ 到 } U_{i})\)）。

